\section{Nums}
Star-CCM+ uses the one-fluid formulation to implement the volume-of-fluid method. 
 The two-fluid model is implemented as the eulerian multiphase (EMP) model.
 The eulerian
 multiphase model is further equipped to resolve large scale interfaces and
 thus can model both disperse and separated flow regimes.~\cite{gadaLargeScaleInterface2017}~\citelater{starccm} 
 This chapter details the numerical implementation of both the models in Star-CCM+ and discusses 
 the closures relations used to close the system of equations.

\subsection{VOF}
The volume-of-fluid (VOF) method is an interface capturing method that utilises the governing 
 equations of the one-fluid model as detailed in Section~\ref{sec:one-fluid}. The method presumes that
 the computational mesh is fine enough to resolve the interface topology.
 Using the volume fraction as the marker function - the interface is assumed to be present
 in cells with volume fraction values between 0 and 1.
 The distribution of the phase $i$ is driven by the volume fraction transport equation,
\begin{equation}\label{eq_vof_transport}
\frac{\partial}{\partial t}\int\limits_{V} \alpha_{i}\,\mathrm{d}V +
\int\limits_{\partial V} \alpha_{i} \vetor{v}_{i} \cdot \mathrm{d}\vetor{a}\;= \;
\int\limits_{V}\left(S^{\alpha}_{i} - \frac{\alpha_{i}}{\rho_{i}} \frac{D \rho_{i}}{D t}\right)\,\mathrm{d}V
- \int\limits_{V} \frac{1}{\rho_{i}}\nabla (\alpha_{i}\rho_{i}\vetor{v}_{d,i})\,\mathrm{d}V\,.
\end{equation}
 In case of two phases, the volume fraction transport is solved only for the first phase, thus in each volume
 element, the volume fraction of phase $j$ is given as $\alpha_{j} = 1 - \alpha_{i}$.
 
The High Resolution Interface Capturing (HRIC) scheme is the default volume fraction convection scheme used.
 The scheme is based on the normalised variable diagrams and blends the downwind
 scheme with the upwind scheme to resolve the interface parallel or perpendicular to the flow direction.
 The calculated normalised face values of the volume fractions are further corrected in order 
 to conform to the local Courant number. The spurious oscillations that may arise in the volume fraction field due to meshes with large aspect ratios
 are suppressed by artificially smoothening the volume fraction gradient across the interface.~\citelater{star-ccm+}
~\COMMENT{how is turbulence implemented?}

\subsection{EMP}
Star-CCM+ implements the governing equations of the two-fluid model as detailed in Section~\ref{sec:two-fluid} as the eulerian multiphase (EMP) model. 
 The two-fluid model is predominantly used to simulate dispersed flows and the mesh resolution is
 assumed to be coarse enough to not capture the interface topology between the phases. The EMP model,
 however, utilises the multiple flow regime framework to allow the simulation of multiscale flows.
 Based on increasing volume fraction of the
 secondary phase, $\alpha_{s}$, in a finite volume cell,
 the topology transitions through three principal 
 regimes~\cite{gadaLargeScaleInterface2017}~\citelater{starccm}, where the secondary phase is:
 \begin{itemize}
    \item dispersed in the primary phase (first dispersed regime, $0 \le \alpha_{s} \le \alpha_{fr} $)
    \item separated from the primary phase through a large scale interface (interfacial flow regime, $\alpha_{fr} \le \alpha_{s} \le \alpha_{sr} $) 
    \item continuous while having the primary phase dispersed (second dispersed regime, $\alpha_{sr} \le \alpha_{s} \le 1 $)
 \end{itemize}
The first and second regime terminuses, $\alpha_{fr}$ and $\alpha_{sr}$, mark the values of $\alpha_{s}$
that defines the end of the first dispersed and large scale interface regimes respectively. Star-CCM+ uses a regime-specific closure models for 
 the interfacial momentum and energy terms. These models are blended using a weighting
 scheme according to the calculated proportion of each regime in the cell.

A variety of convection schemes for the transport of volume fraction are available. The Adaptive Interface Sharpening (ADIS) scheme can be
 used in the multiple flow regime framework to avoid spurious oscillations near sharp interfaces due to
 the upwind schemes,
 or artificial sharpening of otherwise smooth regions due to HRIC.
 ADIS uses the HRIC scheme in the vicinity of a large scale interface and uses a standard 
 upwind scheme away from the sharp interface.~\citelater{starccm}
 %turbulence damping near interfaces?
~\COMMENT{how is turbulence implemented?}

\subsection{Star-CCM+ Stuff - needs a better heading}
~\COMMENT{relevance?}
%multiple flow regimes - blinkov 1993 plots a simple regime map where bubbly flow is \alpha<0.3 and droplet flow at \alpha>0.7
The relevant variables are
 calculated at the center of each control volume and interpolated to the control volume surfaces.
 Convective and diffusive fluxes are interpolated using the second-order upwind differentiation
 scheme which uses the linear interpolation between two upwind centroids to approximate
 the fluxes at the surfaces of the control volume. Gradients are interpolated using 
 the least squares method with Venkatakrishnan gradient limiters.
 The semi-implicit method for pressure-linked equations (SIMPLE) algorithm is used 
 as the default segregated pressure-velocity coupling scheme. 
~\COMMENT{what to mention later: volume fraction convection, semi-implicit treatment, blending schemes and stuff.}


The implementation of the SST k-$\omega$ model in Star-CCM+ can be further controlled to
 adjust the realisability. Realisability is the positive semi-definiteness of the 
 reynolds stress tensor. The specification of turbulence is done using
 the turbulence intensity and the length scale. The inherent anisotropy 
 in turbulence is accounted for with quadratic and cubic constitutive options
 in the modelling of reynolds stresses. The parameters in the model
 are left on their default values as they are empirically modelled.

 At large scale interfaces, the near-interface flow often exhibits
 laminar characteristics with the development of boundary layers. However, in situations
 with significant velocity gradients, turbulent fluctuations can arise in the
 regions adjacent to the interface (especially when the interface is mobile or
 highly distorted). Recognizing that standard RANS turbulence
 models tend to overpredict turbulence intensity near phase boundaries,
 STAR-CCM+ allows the user to apply turbulence damping strategies
 (e.g., vorticity limiters) which suppress excessive eddy viscosity and
 restore physically consistent laminar sublayers adjacent to the interface.
%"neither of the two methodology - VOF or two-fluid - is suitable for multiscale flows~\cite{gadaLargeScaleInterface2017}".
%"most data for the semi-empirical relations covers rather narrow domains~\cite{yadigarogluNatureMultiphaseFlows2018}"
%"if one considers a non flow situation in the two horizontal layers, the forces acting on the 
%two fluids and their interface are gravity and surface tension - in this case we speaf of
%the RT instability~\cite{hewittModellingStrategiesTwoPhase2018}."
%"there exists a wide spectrum of scales, ranging from the global one dictated by the flow
%configuration, to the smalles one where turbulent energy is dissipated into heat\cite{shyyComputationalTechniquesComplex1997}"

%While there are multiple numerical methods to account for volume fraction transport, the 
% High-Resolution Interface Capturing scheme is specifically implemented to handle 
% sharp phase interfaces. Free-surface flows with sharp 
% interfaces present a discontinuity
% in the solution domain, causing discretisation errors. This can lead to
% spurious or parasitic currents.~\cite{MIRJALILI2019221}
% Although often negligible, these currents can be accommodated for through 
% the incorporation of an artificial viscosity term.   


\subsection{Closures}\label{sec:closures}
Closures are critical to the accurate prediction of the relevant flow variables and the distribution
 of volume fraction~\cite{sugrueRobustMomentumClosure}. Existing literature reviews 
 the overwhelming number of closures available that yield a range of varyinh performance 
 along with their comparison with experimental data, however literature
 presenting CFD simulations tend to choose 'result-oriented' closures that aren't
 transferable to a wider range of flow situations~\cite{krepperValidationClosureModel2018, sugrueAssessmentSimplifiedSet2017}.
 Literature admits that the industry relies on the tunable parameters in the closures to 
 match the specific experimental data being simulated or the arbitrariness of this choice
 to get improved flow predictions~\citelater{[20] rzehak, kriebitzsch; silva [21] (sugrue)}.
 
Models may also provide satisfying predictions for certain parameters or situations while
 failing in others and a general inconsenus prevails while discussing the effects of interfacial
 forces. Early work has tended to neglect 
 the effects of lift or virtual mass forces, or chosen to use constant coefficients to 
 describe them while using formulations for the drag forces based on 
 relevant non-dimensional numbers.~\cite{liaoEulerianModellingTurbulent2018}
 The effect of non-drag forces, however, tends to be as significant on the 
 distribution of phases.~\citelater{liao2014-flashboiling}
 Since it is 
 difficult to separate the effects of the various forces in actual experiments,
 the development of these closures assumes certain idealisations when they're not
 partly or wholly based on empirical data, rendering them highly flow 
 dependent~\citelater{troshko and hassan, [15] in Sugrue}. The authors 
 of~\citelater{troshko and hassan} critiques the methodology of resolving a 
 complex interfacial momentum transfer through a linear superposition 
 of several simple forces as the basis of the discrepancy.
 
The industry has, however, attempted to establish generality that can be transferred 
 to multiple flow situations and has proposed baseline frameworks. There is a general consensus
 that the role of turbulence is underestimated~\cite{colomboBenchmarkingCFDModelling};
 current models also aim to reduce 
 uncertainty and inconsistency in model selection~\cite{colomboBenchmarkingCFDModelling}, 
 adopt closures that avoid "case-specific adjustments to unphysical coefficients and 
 tunable parameters"~\cite{sugrueRobustMomentumClosure}, and utilise more
 physics based closures that are more physically consistent~\cite{liaoEulerianModellingTurbulent2018,colomboBenchmarkingCFDModelling}.

Considering the 

Drag force is generally more dominant than other forces as they tend to be a magnitude
 higher than other forces~\cite{tas-koehlerCriticalAnalysisDrag2021a}. Most drag models perform 
 equally well at low gas volume fractions. However as they are based on ideal conditions, 
 they do not account for bubble swarm effects at higher gas volume fractions and show reduced 
 accuracy. Certain swarm corrections do exist, but they are averaged along a radial profile when
 the swarm effects are also felt in the bulk~\cite{sugrueRobustMomentumClosure, sugrueAssessmentSimplifiedSet2017,colomboBenchmarkingCFDModelling}.
 The accurate modelling of interaction area density and length becomes relevant for the calculation
 of the drag force. 
 Using the linearised drag coefficient, $A_{D}$, and the relative velocity between the phases
 $\vetor{v}_{r} = \vetor{v}_{j} - \vetor{v}_{i}$, Star-CCM+ implements the drag force as,
 \begin{equation}
    \vetor{F}_{ij}^{D} = A_{D}\vetor{v}_{r}
 \end{equation}
 The linearised drag coefficient is related to the standard definitions of the drag 
 coefficient $C_{D}$ using
 \begin{equation}
    A_{D} = C_{D} \frac{1}{2} \rho_{c} |\vetor{v}_{r}| \frac{a_{cd}}{4}
 \end{equation}
 where $a_{cd}$ is the interfacial area density and $a_{cd}/4$ represents the projected area of
 the equivalent spherical particle. 

The lift force models are particularly sensitive and uncertain in their performance~\cite{sugrueAssessmentSimplifiedSet2017}.
 They suffer from the same discrepancy between performance in low and 
 high gas volume fractions. Moreover, existing 
 lift force implementations in Star-CCM+ fail to replicate the experimental data
 they were derived from~\cite{sugrueRobustMomentumClosure}.
 Assuming $\alpha_{p}$ is the volume fraction of the primary (continuous)
 phase, Star-CCM+ implements the lift force as
 \begin{equation}
    F_{ij}^{L} = C_{L,\text{eff}} \alpha_{d} \rho_{p} 
    \left[\vetor{v}_{r} \times \left(\nabla \times \vetor{v}_{p}\right)\right].
 \end{equation} 

A turbulent dispersion force is supposed to act on the gas phase and 
 spreads this dispersed phase due to turbulence
 in the continuous phase. This is typically modelled as 
 as a diffusion term proportional to the turbulent 
 kinetic energy gradients~\cite{liaoFlashingEvaporationDifferent2013}. 
 The force is theorised to arise from the averaging of the fluctuations
 of the bubbles due to the turbulence.
 \begin{equation}
    F_{ij}^{TD} = A_{D} \vetor{v}_{TD}
    = A_{D}\,\underbrace{C_{0} \frac{\nu_{c}^{t}}{\sigma_{\alpha}}\tens{I}}_{\tens{D}_{TD}} \cdot 
    \left[
        \nabla \ln(\alpha_{d}) - \nabla \ln(\alpha_{c})
        \right]
 \end{equation}
 The closure relation by Burns~\citelater{burns} uses the Favre averaging of the drag
 force and can be expressed as
 \begin{equation}
    F_{ij}^{TD} = C_{D} C_{TD} \frac{\mu_{l}^{t}}{\rho_{c} \sigma_{TD}}\vetor{v}_{TD},
 \end{equation}
 where $C_{D}$ is the drag coefficient, $C_{TD}$ is an adjustable turbulent dispersion
 calibration coefficient, $\mu_{l}^{t}$ is the liquid turbulent viscosity, and
 $\sigma_{TD}$ is the turbulent Prandtl number.

Most lift
 coefficient correlations don't account for the decrease of the magnitude of the
 lift force to zero 
 near the walls. This low gas volume fraction at the walls has historically
 been explained
 through a liquid film forming at the walls repelling the bubbles. This has been demonstrated to
 not actually occur in practice. The wall lubrication force has thus been criticised
 as lacking any physical justification~\cite{sugrueRobustMomentumClosure,colomboBenchmarkingCFDModelling}.
 Newer closures~\citelater{lubchenko} utilise a turbulent dispersion force based 
 formualtion to maintain the desired gas volume fraction and use the closure instead
 as a correction term to resolve the phase distribution at the boundary lost due to
 the averaging process. Star-CCM+'s implementation of the wall lubrication force
 \begin{equation}
    F_{ij}^{WL} = - C_{WL}(y_{w})\alpha_{s} \rho_{p} 
    \frac{|\vetor{v}_{r} - (\vetor{v}\cdot \vetor{n})\vetor{n}|}{d_{b}} \vetor{n}
 \end{equation}
 can be changed to implement alternative models instead by defining user defined 
 field functions for $C_{WL}$.

Surface tension, although an interfacial force, is implemented using the 
 Continuum Surface Force approach as a volumetric force in the VOF formulation.
 Divided into normal and tangential components and using the mean 
 curvature of the free surface 
 $\kappa = - \nabla \cdot \left(\nabla \alpha_{i}/|\nabla \alpha_{i}|\right)$,
 the volumetric surface tension force is given as 
 \begin{equation}
    \vetor{f}_{\sigma} = \underbrace{-\sigma \nabla \cdot \frac{\nabla \alpha_{i}}{|\nabla \alpha_{i}|} \nabla \alpha_{i}}_{\vetor{f}_{\sigma,n}}
     + \underbrace{(\nabla \sigma)_{t}|\nabla \alpha_{i}|}_{\vetor{f}_{\sigma,t}}
 \end{equation}
 The EMP treats surface tension by splitting it among the phases occupying
 the cell using the local volume fractions. The pressure gradient
 is then the summation of the 
 individual phase surface tension contributions $\vetor{f}_{\sigma, i}$~\cite{gadaLargeScaleInterface2017}~\citelater{starccm},
 \begin{equation}
    \vetor{f}_{\sigma, i} = \alpha_{i} \sigma \kappa \nabla \alpha_{i}.
 \end{equation}

\subsection{Closures: Under Construction}
\begin{itemize}
    \item Interphase mass transfer (evaporation/condensation)
     Represents phase change, modeled as interfacial flux determined
by local thermodynamic properties (temperature, concentration, pressure—often
using kinetic models for flash evaporation).
    \item Interphase energy transfer
     Encompasses latent heat and sensible heat exchange at the interface 
and throughout the phases (essential for accurate phase change modeling).
- conduction component based on the 'thin thermal boundary layer' assumption
    \item Nucleation (wall and bulk)
    \item Interfacial area changes due to bubble growth, coalescence, breakup
    \item Local pressure and temperature gradients driving simultaneous mass and energy transfer
\end{itemize}

\subsection{Vague and Miscellaneous}